\documentclass{article}
\usepackage{amsmath}
\usepackage{mathtools}
\usepackage{amssymb}
\begin{document}

\DeclarePairedDelimiter\bra{\langle}{\rvert}
\DeclarePairedDelimiter\ket{\lvert}{\rangle}
\DeclarePairedDelimiterX\braket[2]{\langle}{\rangle}{#1\,\delimsize\vert\,\mathopen{}#2}

Dans cette partie, nous allons utiliser une solution aux équations de Maxwell qui va être très
importante.
Un rappel sur les équations de Maxwell \\
Loi de Gauss pour l'électricité:
$$
\nabla \cdot \vec{E} = \frac{\rho}{\epsilon_0}
$$
Loi de Gauss pour le magnétisme:
$$
\nabla \cdot \vec{B} = 0
$$
Loi de Faraday pour l'induction:
$$
\nabla \times \vec{E} = -\frac{\partial\vec{B}}{\partial t}
$$
Loi d'Ampère-Maxwell:
$$
\nabla \times \vec{B} = \mu_0 J + \epsilon_0 \mu_0 \frac{\partial\vec{E}}{\partial t}
$$
Une solution très utile pour ces équations est la suivante:
$$
\vec{E} = e^{i(kz-\omega t)}(E_x \hat{i} + E_y \hat{j})
$$
$$
\vec{B} = \frac{1}{\omega} \hat{k} \times \vec{E}
$$
V est le potentiel électrique défini comme étant:
$$
\vec{E} = - \nabla V = (-\frac{\partial V}{\partial x},-\frac{\partial V}{\partial y},-\frac{\partial V}{\partial z})
$$
Le curl \( \nabla \times \vec{F} \) d'un champ vectoriel est défini comme étant:
$$
\nabla \times \vec{F} = 
\begin{vmatrix}
	\hat{i} & \hat{j} & \hat{k} \\
	\frac{\partial}{\partial x} & \frac{\partial}{\partial y} & \frac{\partial}{\partial z} \\
	F_x & F_y & F_z
\end{vmatrix}
$$
Pour trouver le vecteur de Jones, nous allons d'abord voir de quels composants est véritablement composé le champ électrique, ensuite nous allons trouver exactement les facteurs dans chaque direction.
$$
\vec{E} = E_x \hat{i} e^{i(kz-\omega t)} + E_y \hat{j} e^{i(kz-\omega t)} 
$$
$$
\vec{B} = \frac{1}{\omega}\,\hat{k} \times \vec{E}
$$
On calcule:
$$
\nabla \times \vec{E} = 
\begin{vmatrix}
	\hat{i} & \hat{j} & \hat{k} \\
	\frac{\partial}{\partial x} & \frac{\partial}{\partial y} & \frac{\partial}{\partial z} \\
	E_x & E_y & E_z
\end{vmatrix}
$$
$$
\iff \nabla \times \vec{E} =
\begin{vmatrix}
	\hat{i} & \hat{j} & \hat{k} \\
	\frac{\partial E}{\partial x} & \frac{\partial E}{\partial y} & \frac{\partial E}{\partial z} \\
	\frac{\partial V}{\partial x} & \frac{\partial V}{\partial y} & \frac{\partial V}{\partial z} \\
\end{vmatrix}
$$
$$
\iff \nabla \times \vec{E} = \hat{i}(\frac{\partial}{\partial y} \cdot -\frac{\partial V}{\partial z} - \frac{\partial}{\partial z} \cdot -\frac{\partial V}{\partial y}) \newline
- \hat{j}(\frac{\partial}{\partial x} \cdot -\frac{\partial V}{\partial z} - \frac{\partial}{\partial z} \cdot -\frac{\partial V}{\partial x}) \newline
+ \hat{k}(\frac{\partial}{\partial x} \cdot -\frac{\partial V}{\partial y} - \frac{\partial}{\partial y} \cdot -\frac{\partial V}{\partial x})
$$
$$
E_z = 0 \therefore \frac{\partial}{\partial z} = \frac{\partial V}{\partial z} = 0
$$
$$
\therefore \nabla \times \vec{E} = \hat{i}(0+0) - \hat{j}(0+0)
+ \hat{k}(\frac{\partial}{\partial x} \cdot -\frac{\partial V}{\partial y} - \frac{\partial}{\partial y} \cdot -\frac{\partial V}{\partial x})
$$
$$
\iff \nabla \times \vec{E} = \hat{k}(\frac{\partial}{\partial x} \cdot -\frac{\partial V}{\partial y} - \frac{\partial}{\partial y} \cdot -\frac{\partial V}{\partial x})
$$
$$
\therefore \nabla \times \vec{E} \perp \vec{E}
$$
Avec la projection de \( E_x \) et de \( E_y \), on peut donc trouver que:
$$
E = E_0 \cos(\theta) \hat{i} e^{i(kz-\omega t)} + E_0 \sin(\theta) \hat{j}
e^{i(kz-\omega t)} \\
$$
$$
\iff E = E_0 (\cos(\theta)\hat{i} + \sin(\theta)\hat{j}) e^{i(kz-\omega t)} \\
$$
$$
\iff E = E_0 e^{i(kz-\omega t)} \ket{\psi}
$$
$$
\iff \ket{\psi} = \begin{pmatrix} \cos(\theta) \\ \sin(\theta) \end{pmatrix}
$$
\( \ket{\psi} \) est donc le vecteur de Jones représentant la direction du champ électrique dans
une polarisation linéaire. Il est composé seulent de \( \hat{i} \) et de \( \hat{j} \),
mais pas de \( \hat{k} \), car il est perpendiculaire à \( \nabla \times {E} \) qui est composé
seulement de \( \hat{k} \).

\end{document}
